\section*{Bab \ref{ch:g9:metals-acids-corrosion} --- Logam: Reaksi dengan Asam dan Korosi}

\begin{solution}{exo:g9:metals-acids-corrosion:1}
Hidrogen: bilah kayu menyala di mulut tabung menimbulkan letupan
melengking.
\end{solution}

\begin{solution}{exo:g9:metals-acids-corrosion:2}
Besi, seng, dan aluminium bereaksi; tembaga dan emas tidak.
\end{solution}

\begin{solution}{exo:g9:metals-acids-corrosion:3}
Oksigen dan air, kedua-duanya.
\end{solution}

\begin{solution}{exo:g9:metals-acids-corrosion:4}
\omterm{def:g9:ions:ion}{Ion} yang ada selama reaksi tetapi tidak ikut bereaksi. Dengan asam
klorida: \omterm{def:g9:ions:ion}{ion} klorida \ce{Cl-}.
\end{solution}

\begin{solution}{exo:g9:metals-acids-corrosion:5}
\ce{Mg(s) + 2H+(aq) -> Mg^{2+}(aq) + H2(g)}; muatan $+2$ di setiap ruas.
\end{solution}

\begin{solution}{exo:g9:metals-acids-corrosion:6}
Tambahkan natrium hidroksida: \omterm{def:g9:ions:precipitate}{endapan} hijau \ce{Fe(OH)2} menunjukkan
\omterm{def:g9:ions:ion}{ion} besi(II).
\end{solution}

\begin{solution}{exo:g9:metals-acids-corrosion:7}
Pendidihan menghilangkan \omterm{def:g4:air-a-mixture-of-gases:air}{udara} yang terlarut dalam air, dan minyak
mencegah \omterm{def:g4:air-a-mixture-of-gases:air}{udara} \omterm{def:g2:mixing-and-dissolving:dissolve}{larut} kembali: paku itu mendapat air tetapi tidak
mendapat oksigen.
\end{solution}

\begin{solution}{exo:g9:metals-acids-corrosion:8}
Garam yang \omterm{def:g2:mixing-and-dissolving:dissolve}{larut} dalam air yang terpercik ke mobil mempercepat
perkaratan.
\end{solution}

\begin{solution}{exo:g9:metals-acids-corrosion:9}
Lapisan oksida yang terbentuk melekat pada logam dan menutupinya dari
\omterm{def:g4:air-a-mixture-of-gases:air}{udara}: serangan itu berhenti di permukaan.
\end{solution}

\begin{solution}{exo:g9:metals-acids-corrosion:10}
Catlah, beri minyak atau gemuk, galvanisasilah (atau buatlah dari baja
tahan \omterm{def:g9:metals-acids-corrosion:corrosion}{karat}).
\end{solution}

\begin{solution}{exo:g9:metals-acids-corrosion:11}
\ce{2Al(s) + 6H+(aq) -> 2Al^{3+}(aq) + 3H2(g)}: 2 Al dan 6 H di setiap
ruas, muatan $+6$ di setiap ruas. Untuk 200 \omterm{def:g7:atoms-and-molecules:atom}{atom} Al:
$200 \times \frac{3}{2} = 300$ \omterm{def:g7:atoms-and-molecules:molecule}{molekul} hidrogen.
\end{solution}

\begin{solution}{exo:g9:metals-acids-corrosion:12}
Seng di sekitar goresan terkorosi menggantikan baja: selama masih ada
seng di dekatnya, bajanya terhindar.
\end{solution}

\begin{solution}{pb:g9:metals-acids-corrosion:1}
\textbf{1.} \ce{Zn(s) + 2H+(aq) -> Zn^{2+}(aq) + H2(g)}.

\textbf{2.} Hidrogen: letupan melengking dengan bilah kayu menyala.

\textbf{3.} Natrium hidroksida memberikan \omterm{def:g9:ions:precipitate}{endapan} putih
\ce{Zn(OH)2}.

\textbf{4.} \omterm{def:g9:ions:ion}{Ion-ion} klorida \ce{Cl-}.

\textbf{5.} $125.6 - 113.5 = \qty{12.1}{g}$.

\textbf{6.} Hidrogen dilepaskan selama seng bereaksi dengan asam; ketika
seng sudah habis, gelembung-gelembungnya berhenti.

\textbf{7.} Besi dalam baja juga akan bereaksi, menghasilkan hidrogen dan
\omterm{def:g9:ions:ion}{ion} \ce{Fe^{2+}}; natrium hidroksida lalu akan memberikan \omterm{def:g9:ions:precipitate}{endapan} hijau.

\textbf{8.} Seng diserang menggantikan baja dan menjauhkan air dan \omterm{def:g4:air-a-mixture-of-gases:air}{udara}
dari baja itu.

\textbf{9.} $2 \times 10.0 \times 10.0 = \qty{200}{cm^2}$.

\textbf{10.} $12.1 \div 7.13 \approx \qty{1.70}{cm^3}$.

\textbf{11.} $1.70 \div 200 = \qty{0.0085}{cm}$.

\textbf{12.} $0.0085 \times 10\,000 = \qty{85}{\micro m}$: lapisan seng
itu tebalnya sekitar \qty{85}{\micro m}.
\end{solution}
